% ======================================================================
% APPENDIX B: OPENFLOW 1.3 PROTOCOL TECHNICAL REFERENCE
% ======================================================================
\chapter{OpenFlow 1.3 Protocol Technical Reference}
\label{app:openflow}

\section{OpenFlow 1.3 Common Header Specification}
All OpenFlow 1.3 protocol messages begin with an 8-byte common header. The binary wire format is detailed in Table~\ref{tab:ofp_header_format}.

\begin{table}[htbp]
\centering
\caption{OpenFlow 1.3 Common Header Binary Structure}
\label{tab:ofp_header_format}
\small
\begin{tabularx}{\textwidth}{c c c X}
\toprule
\textbf{Byte Offset} & \textbf{Field Name} & \textbf{Width} & \textbf{Description \& Wire Value} \\
\midrule
0 & \texttt{version} & 8 bits & Protocol version. For OpenFlow 1.3, wire value is \texttt{0x04}. \\
1 & \texttt{type} & 8 bits & Message type code (e.g., \texttt{OFPT\_PACKET\_IN = 10}). \\
2--3 & \texttt{length} & 16 bits & Total length of message including header in bytes. \\
4--7 & \texttt{xid} & 32 bits & Transaction ID associated with this message for matching requests to replies. \\
\bottomrule
\end{tabularx}
\end{table}

\section{Core OpenFlow 1.3 Message Type Enumeration}
Table~\ref{tab:ofp_types_enumeration} lists the primary message types defined in ONF TS-006 \cite{onf2012openflow} utilized by our platform.

\begin{table}[htbp]
\centering
\caption{Primary OpenFlow 1.3 Message Types Utilized in the Platform}
\label{tab:ofp_types_enumeration}
\small
\begin{tabularx}{\textwidth}{c l c X}
\toprule
\textbf{Code} & \textbf{Message Constant} & \textbf{Direction} & \textbf{Operational Function} \\
\midrule
0 & \texttt{OFPT\_HELLO} & Symmetric & Negotiates protocol version during TCP connection setup. \\
2 & \texttt{OFPT\_ECHO\_REQUEST} & Symmetric & Probes connection liveness and active Echo-RTT latency. \\
3 & \texttt{OFPT\_ECHO\_REPLY} & Symmetric & Responds to echo request with identical payload. \\
5 & \texttt{OFPT\_FEATURES\_REQUEST} & Controller $\to$ Switch & Queries switch DPID, buffer count, and table capabilities. \\
6 & \texttt{OFPT\_FEATURES\_REPLY} & Switch $\to$ Controller & Reports datapath ID and supported OpenFlow capabilities. \\
10 & \texttt{OFPT\_PACKET\_IN} & Switch $\to$ Controller & Transfers unclassified packet header/miss to controller. \\
13 & \texttt{OFPT\_PACKET\_OUT} & Controller $\to$ Switch & Instructs switch to output packet out of specific port(s). \\
14 & \texttt{OFPT\_FLOW\_MOD} & Controller $\to$ Switch & Installs, modifies, or deletes flow table entries in TCAM. \\
15 & \texttt{OFPT\_GROUP\_MOD} & Controller $\to$ Switch & Manages OpenFlow Group tables (Select / Fast Failover). \\
18 & \texttt{OFPT\_MULTIPART\_REQUEST} & Controller $\to$ Switch & Polls port stats, flow counters, or table descriptions. \\
19 & \texttt{OFPT\_MULTIPART\_REPLY} & Switch $\to$ Controller & Delivers multi-part port octets and packet counters. \\
20 & \texttt{OFPT\_BARRIER\_REQUEST} & Controller $\to$ Switch & Enforces sequential execution of preceding FlowMod rules. \\
21 & \texttt{OFPT\_BARRIER\_REPLY} & Switch $\to$ Controller & Confirms all prior FlowMod rules are committed to TCAM. \\
\bottomrule
\end{tabularx}
\end{table}

\section{Annotated OpenFlow Hex Packet Dumps}

\subsection{Annotated \texttt{OFPPacketIn} Wire Frame}
Listing~\ref{lst:packet_in_hex} displays an annotated hex dump of an \texttt{OFPPacketIn} message captured during live Mininet emulation when a new flow miss triggers an event.

\begin{lstlisting}[basicstyle=\ttfamily\scriptsize, caption={Annotated Hex Dump of OFPPacketIn Message}, label={lst:packet_in_hex}]
04 0a 00 52 00 00 00 07   # Version: 0x04 (OF 1.3), Type: 0x0a (PACKET_IN), Len: 82 bytes, XID: 7
00 00 00 00 00 00 00 00   # Buffer ID: 0x00000000 (No Buffer, full packet encapsulated)
00 3c 01 00 00 00 00 00   # Total Len: 60 bytes, Reason: 0x01 (OFPR_ACTION_MISS), Table ID: 0
00 00 00 00 00 00 00 01   # Cookie: 0x0000000000000001
00 01 00 0c 80 00 00 04   # Match: Type 1 (OXM), Length 12 bytes, OXM_OF_IN_PORT header
00 00 00 02 00 00 00 00   # Ingress Port: Port 2, 4-byte padding to 8-byte alignment
# --- Encapsulated Ethernet Frame (IPv4 / TCP SYN) ---
00 00 00 00 00 02 00 00   # Dest MAC: 00:00:00:00:00:02 (Host h2)
00 00 00 01 08 00 45 00   # Src MAC:  00:00:00:00:00:01 (Host h1), EtherType: 0x0800 (IPv4)
00 28 00 01 00 00 40 06   # IPv4 Header: Len 40, TTL 64, Protocol: 0x06 (TCP)
7c c0 0a 00 00 01 0a 00   # Checksum, Src IP: 10.0.0.1 (h1)
00 02 1f 90 00 50 00 00   # Dst IP: 10.0.0.2 (h2), Src Port: 8080, Dst Port: 80 (HTTP)
\end{lstlisting}

\subsection{Annotated \texttt{OFPFlowMod} Installation Wire Frame}
Listing~\ref{lst:flow_mod_hex} displays the corresponding \texttt{OFPFlowMod} message generated by \texttt{RoutingModule} pushing a Priority 10 forwarding rule.

\begin{lstlisting}[basicstyle=\ttfamily\scriptsize, caption={Annotated Hex Dump of OFPFlowMod Installation Message}, label={lst:flow_mod_hex}]
04 0e 00 68 00 00 00 08   # Version: 0x04 (OF 1.3), Type: 0x0e (FLOW_MOD), Len: 104 bytes, XID: 8
00 00 00 00 00 00 00 00   # Cookie: 0x0, Cookie Mask: 0x0
00 00 00 00 00 00 00 00   # Table ID: 0, Command: 0x00 (OFPFC_ADD), Idle Timeout: 60s
00 3c 00 78 00 0a 00 00   # Hard Timeout: 120s, Priority: 10 (0x000a), Buffer ID: 0xffffffff
# --- Match Fields (OXM) ---
00 01 00 1a 80 00 0a 02   # OXM Type 1, Length 26, EtherType Match Header
08 00 80 00 16 04 0a 00   # EtherType: 0x0800, IPv4 Src Match: 10.0.0.1
00 01 80 00 18 04 0a 00   # IPv4 Dst Match: 10.0.0.2
00 02 00 00 00 00 00 00   # Padding to 8-byte alignment
# --- Instructions (APPLY_ACTIONS: Output to Port 3) ---
00 04 00 18 00 00 00 00   # Instruction: 0x0004 (OFPIT_APPLY_ACTIONS), Length 24 bytes
00 00 00 10 00 00 00 03   # Action: 0x0000 (OFPAT_OUTPUT), Length 16, Output Port: 3
00 00 00 00 00 00 00 00   # Max Length: 0 (Wire speed), Padding
\end{lstlisting}
